\chapter{Conditions}

\section{Clostridium difficile infection (c. diff)}
\label{cond:c-diff}

\section{Peptic ulcer}
\label{cond:peptic_ulcer}

Begins with a break in the epithelium which can result from e.g. \hyperref[cond:h-pylori]{h. pylori}, \hyperref[drugclass:NSAIDs]{NSAIDs}

30\% of ulcers are gastric, and 70\& are duedenal \cite[7:32]{CobboldPepticUlcerDisease}.

\section{Helicobacter pylori (h. pylori)}
\label{cond:h-pylori}

Often acquired through oral-oral or oral-fecal in children \cite{CobboldPepticUlcerDisease}

h. pylori can active G cells, which cause the release of more gastrin, and then more gastric acid, which results in mucosal injury \cite[t. 5:50]{CobboldPepticUlcerDisease}. 

Gastrin also causes a release of pepsinogen, which turns to pepsin under the influence of acid, which then digests exposed mucosal proteins \cite[t. 6:47]{CobboldPepticUlcerDisease}. Takeaway message is that pepsin contributes to ulcer deepening, not the cause. 

\section{Gastro-oesophageal reflux disease (GORD/GERD)}

GORD is any continuous or regular reflux. Anything that disrupts the lower-oesophageal sphincter can lead to GORD

Long term GORD can cause reflux oesophagitus (mucosal redness and itching). Eventually you can get bleeding and scarring in the oseophagus, and then extremely you can get oesophageal narrowing due to scarring and finally Barrett's oesophagus. \cite[9:--]{CobboldGastroOesophagealRefluxDisease}

\textbf{Causes}

\begin{itemize}
    \item Alcohol and caffeine can cause lower-oesophageal relaxation
\end{itemize}